% from main.tex
%%%%%%%%%%  scan idx 281  |  folio 268  %%%%%%%%%%
\item Si \(E\) est
parfait, les \(H^i(E_x\otimes_{\mathcal O_{X,x}}^{\mathrm L}k(x))\), pour \(x\in\pt(X)\), sont des espaces vectoriels de rang fini sur \(k(x)\), nuls sauf un nombre fini, et la fonction
\[
 x\longmapsto\sum(-1)^i\operatorname{rg}_{k(x)}H^i(E_x\otimes_{\mathcal O_{X,x}}^{\mathrm L}k(x))
\]
est localement constante sur \(\pt(X)\).
\item Si \(E\) est \((p-1)\)-pseudo-cohérent, \(H^p(E_x\otimes_{\mathcal O_{X,x}}^{\mathrm L}k(x))\) est, pour \(x\in\pt(X)\), un espace vectoriel de rang fini sur \(k(x)\), et la fonction
\[
 x\longmapsto\operatorname{rg}_{k(x)}H^p(E_x\otimes_{\mathcal O_{X,x}}^{\mathrm L}k(x))
\]
est semi-continue supérieurement sur \(\pt(X)\). Si en outre le foncteur \(M\mapsto H^p(E\otimes^{\mathrm L}M)\) de \(\Mod(X)\) dans \(\Mod(X)\) est exact, \(H^p(E)\) est localement libre de type fini et la fonction précédente est localement constante.
\item Supposons \(E\) \(p\)-pseudo-cohérent (resp. \((p-1)\)-pseudo-cohérent), et soit \(x\in\pt(X)\). Si le foncteur \(M\mapsto H^p(E_x\otimes_{\mathcal O_{X,x}}^{\mathrm L}M)\) est exact à droite (resp. à gauche) sur la catégorie \(\Mod(x)\) des \(\mathcal O_{X,x}\)-Modules, il existe un ouvert \(x\)-ponctué \(U\) de \(X\) tel que le foncteur \(M\mapsto H^p(E\,|\,U\otimes^{\mathrm L}M)\) soit exact à droite (resp. à gauche) sur \(\Mod(U)\).
\end{enumerate}

\emph{Preuve.} (i) D'après (I 4.19.1), \(E_x\otimes_{\mathcal O_{X,x}}^{\mathrm L}k(x)\) est un complexe parfait de \(k(x)\)-modules, donc (I 5.9 et 3.5) les \(H^i(E_x\otimes_{\mathcal O_{X,x}}^{\mathrm L}k(x))\) sont de rang fini sur \(k(x)\) et nuls sauf un nombre fini. On a\footnote{\(\operatorname{rg}(\ )\) désigne le rang d'un complexe parfait (I 6.11).}
\begin{equation}
 \operatorname{rg}(E_x\otimes_{\mathcal O_{X,x}}^{\mathrm L}k(x))=\sum(-1)^i\operatorname{rg}_{k(x)}H^i(E_x\otimes_{\mathcal O_{X,x}}^{\mathrm L}k(x)),\tag{5.5.1}
\end{equation}
comme on le voit par un dévissage immédiat utilisant l'additivité de la fonction rang. Compte tenu de (5.4.1 a)), l'assertion (i) résulte donc de la
